% Generated draft: requires mathematical and source review.
\documentclass[11pt]{article}
\usepackage[margin=25mm]{geometry}
\usepackage{fontspec}
\setmainfont{Latin Modern Roman}
\usepackage{amsmath,amssymb,mathtools,unicode-math}
\setmathfont{Latin Modern Math}
\usepackage{enumitem,needspace,xurl,hyperref}
\setlength{\parindent}{0pt}
\setlength{\parskip}{5pt}
\setlength{\emergencystretch}{3em}
\newcommand{\N}{\mathbb{N}}
\newcommand{\Z}{\mathbb{Z}}
\newcommand{\Q}{\mathbb{Q}}
\newcommand{\R}{\mathbb{R}}
\newcommand{\C}{\mathbb{C}}
\newcommand{\F}{\mathbb{F}}
\newcommand{\A}{\mathbb{A}}
\newcommand{\PP}{\mathbb{P}}
\newcommand{\E}{\mathbb{E}}
\newcommand{\eps}{\varepsilon}
\newcommand{\dd}{\,\mathrm{d}}
\begin{document}
\section*{UnsolvedMath ID: 700709}
\textbf{Source problem code:} catalogue entry
\subsection{Definitions and mathematical statement}
Let \(r\) be a positive integer. This is the range supported by the supplied recursive notation; the source does not explicitly state the range of \(r\) or define a zeroth iterated commutator. For this draft, use the convention \([x,y]=x^{-1}y^{-1}xy\), and define
\[
[x,{}_1 y]=[x,y],\qquad [x,{}_{i+1}y]=[[x,{}_i y],y]\quad(i\geq 1).
\]
The source does not specify its convention for the basic commutator, so this convention requires confirmation against the original source. Define
\[
V_r=\left\{G:\ G\text{ is a finite group and }\exists s\in\Z_{>0},\ s>r,\ \forall x,y\in G,\ [x,{}_r y]=[x,{}_s y]\right\}.
\]
Here \(s=s(G)\) may depend on \(G\), with \(r\) fixed, but must be independent of \(x\) and \(y\).

For a finite soluble group \(G\), its Fitting length \(h(G)\) is the least integer \(n\geq 0\) for which there is a series
\[
1=G_0\trianglelefteq G_1\trianglelefteq\cdots\trianglelefteq G_n=G
\]
with every \(G_i\) normal in \(G\) and every factor \(G_i/G_{i-1}\) nilpotent. In particular, \(h(1)=0\).

\begin{enumerate}
\item[(a)] Does there exist a function \(f:\Z_{>0}\to\Z_{>0}\) such that
\[
\forall r\in\Z_{>0}\ \forall G\in V_r,\qquad G\text{ soluble}\ \Longrightarrow\ h(G)<f(r)?
\]
The integer-valued formulation expresses the requested uniform bound on the integer-valued Fitting length. For \(r<3\), the supplied source directs the reader to the 1983 paper by R. Brandl listed below; it does not state the precise result of that paper.
\item[(b)] For every fixed positive integer \(r\), are there only finitely many isomorphism classes of nonabelian finite simple groups belonging to \(V_r\)? Equivalently, is
\[
\left\{[S]_{\cong}:\ S\in V_r\text{ and }S\text{ is nonabelian simple}\right\}
\]
finite for every such \(r\)? Counting up to isomorphism is the interpretation used here for the source's phrase ``finitely many nonabelian simple groups.'' The supplied source states that this holds for \(r<4\); that assertion is retained as a source qualification and has not been independently verified.
\end{enumerate}
\subsection{Short English statement}
For each fixed positive integer \(r\), consider finite groups whose \(r\)-fold iterated commutator equals a later iterated commutator identically, with the later index allowed to depend on the group. The questions ask for a bound depending only on \(r\) on the Fitting length of soluble members, and for finiteness up to isomorphism of the nonabelian simple members. The supplied source refers to Brandl's work for part (a) when \(r<3\) and reports part (b) as true when \(r<4\).
\subsection{Sources}
\begin{itemize}
\item Supplied local UnsolvedMath registry record, permanent ID \(700709\), exact source problem code \texttt{catalogue entry}, title ``Kourovka Notebook Problem 11.16.'' The registry labels the record \texttt{open}; this is a supplied status label, not an independently verified assessment. The record was created and updated on 2026-10-16 and gives upstream revision \texttt{d371ee917ac9647fbb06cd6fc529deec6f3dda82}.
\item \emph{The Kourovka Notebook (Unsolved Problems in Group Theory)}, original problem 11.16, proposed by R. Brandl, 11th issue (1990). The supplied background also identifies the 21st issue edition (2026), \texttt{arXiv:1401.0300v48}. The supplied collection description attributes editing to E. I. Khukhro and V. D. Mazurov. These bibliographic details are supplied attribution; the cited edition has not been consulted here.
\item R. Brandl, \emph{Bull. Austral. Math. Soc.} \textbf{28} (1983), 101--110. The source cites this reference for part (a) when \(r<3\). No article title, theorem number, or external link is supplied.
\end{itemize}
\end{document}
